\scanpage{106}

La forme plus faible, ``géométrique'' de la
conjecture ci-dessus dirait que \deleted{\unsure{c'est|aussi}} l'on obtient
bien ainsi, par le \unsure{schémoïde|schémaïde} associé à
$M^{\unsure{0|\circ}}\otimes_{\mathbb Z}\mathbb Q$ (sur lequel $\pi$ opère \ldots), la
\deleted{\unsure{opération|grande}} ``partie compacte'' du \deleted{\unsure{sous}} groupe de Galois
motivique de $\overline{\mathbb F}_p$, qui serait donc isom. à

\[
G\simeq D(M_0\otimes\mathbb Q)\times\mathbb G_m
\]

avec \unsure{l'hom|les hom|l'homom.} inj : $\mathbb G_m\longrightarrow G$ s'identifiant
à $j$, et \deleted{$\unsure{(x,\lambda)}$} $(x,\lambda)\unsure{\longmapsto|-}\lambda^{2\deleted{\unsure{x}}}$ s'identifiant à $\varepsilon$.

En termes de valeurs propres de Frobenius :

\uline{Conjecture faible}. Tout $\lambda\in\overline{\mathbb Q}^{*}$, \deleted{\unsure{de la|tel que} tel qu'il}
% slanted
\margin{\unsure{Résultera|Résulterait} \deleted{\unsure{peut}} \unsure{probablement} de la théorie des multiplications complexes : \unsure{cf|cf.} \uline{\unsure{Tanyama-Shimura|Taniyama-Shimura}}}
\deleted{existe un entier $n>0$ avec $\lambda^n$ satisfaisant aux
conditions, à \unsure{savoir} p} un \uline{entier} algébrique,
\deleted{tel qu'il existe un entier $n$}
tel qu'il existe un entier $i\geq0$, de sorte
que $\lambda=p^{i/2}u$, \deleted{soit un} où $u$ est une \uline{unité}
\uline{absolue} \uline{exc}. $p$, et qu'il existe un entier
$n>0$ tel que $u^n\in\mathbb Q'(i)^*$, et $v_p(u^n)+v_{p'}(u^n)=0$
pour \unsure{ts|tt} couples d'idéaux premiers de $\mathbb Q'(i)$ \unsure{induisant|relevant|au-dessus}
\unsure{sur le|sur|d'un} \unsure{$=$|même} idéal premier de $\mathbb Q'$, Alors il
existe un entier $m>0$ tel que
$\lambda^m=(p^m)^{i/2}u^m$ soit \inserted{une} valeur propre de frobenius \unsure{$f_{q^m}$|géom.|$f_{p^m}$}
dans un motif semi-simple \inserted{de poids $i$} convenable sur $\mathbb F_{\unsure{q^m|p^m}}$

% archive sheet number: 105
